\documentclass[11pt]{article}
\usepackage[utf8]{inputenc}
\usepackage[margin=1in]{geometry}
\usepackage{amsmath,amssymb,amsthm}
\usepackage{hyperref}

\newtheorem{theorem}{Theorem}
\newtheorem{lemma}[theorem]{Lemma}
\newtheorem{proposition}[theorem]{Proposition}
\theoremstyle{definition}
\newtheorem{definition}{Definition}
\newtheorem{example}{Example}

\title{Groups, Rings, Fields}
\author{Math Foundations --- Node 04}
\date{}

\begin{document}
\maketitle

\section{Motivation}

Abstract algebra studies sets equipped with operations. By isolating the
\emph{axioms} that an operation satisfies, we can simultaneously reason about
integers, rotations of a cube, residues modulo a prime, invertible matrices,
and continuous symmetries of a manifold. Three structures form the spine of
the subject: groups, rings, and fields, with each obtained by adding more
structure to the previous one.

\section{Groups}

\begin{definition}[Group]
A \emph{group} is a pair $(G, \cdot)$ where $G$ is a set and
$\cdot : G \times G \to G$ is a binary operation satisfying:
\begin{enumerate}
  \item \textbf{Associativity:} $(a \cdot b) \cdot c = a \cdot (b \cdot c)$
        for all $a, b, c \in G$.
  \item \textbf{Identity:} there exists $e \in G$ with
        $e \cdot a = a \cdot e = a$ for every $a \in G$.
  \item \textbf{Inverses:} for every $a \in G$ there exists $a^{-1} \in G$
        with $a \cdot a^{-1} = a^{-1} \cdot a = e$.
\end{enumerate}
A group is \emph{abelian} if additionally $a \cdot b = b \cdot a$ for all
$a, b \in G$.
\end{definition}

\begin{example}
$(\mathbb{Z}, +)$ is an abelian group with identity $0$ and inverses
$n \mapsto -n$. The symmetric group $S_n$ of permutations of $n$ symbols is a
non-abelian group for $n \ge 3$.
\end{example}

\begin{theorem}[Uniqueness of the identity]
\label{thm:identity-unique}
In a group $(G, \cdot)$, the identity element is unique.
\end{theorem}

\begin{proof}
Suppose $e$ and $e'$ are both identity elements of $G$. Applying the identity
axiom for $e'$ to the element $e$, we have $e \cdot e' = e$. Applying the
identity axiom for $e$ to the element $e'$, we have $e \cdot e' = e'$.
Therefore
\[
  e = e \cdot e' = e',
\]
so $e = e'$, and the identity element is unique.
\end{proof}

\begin{proposition}[Uniqueness of inverses]
For each $a \in G$ the inverse $a^{-1}$ is unique.
\end{proposition}

\begin{proof}
Suppose $b$ and $c$ both satisfy $a \cdot b = b \cdot a = e$ and
$a \cdot c = c \cdot a = e$. Then
\[
  b = b \cdot e = b \cdot (a \cdot c) = (b \cdot a) \cdot c = e \cdot c = c,
\]
using associativity (in the middle) and the identity axiom (at the ends).
\end{proof}

\section{Rings}

\begin{definition}[Ring]
A \emph{ring} is a triple $(R, +, \cdot)$ such that
\begin{enumerate}
  \item $(R, +)$ is an abelian group, with identity denoted $0$.
  \item $(R, \cdot)$ is a monoid: $\cdot$ is associative and admits an
        identity $1 \in R$ such that $1 \cdot a = a \cdot 1 = a$ for all
        $a \in R$.
  \item \textbf{Distributivity:} for all $a, b, c \in R$,
        \[
          a \cdot (b + c) = a \cdot b + a \cdot c
          \quad \text{and} \quad
          (b + c) \cdot a = b \cdot a + c \cdot a.
        \]
\end{enumerate}
A ring is \emph{commutative} if $a \cdot b = b \cdot a$ for all $a, b \in R$.
\end{definition}

\begin{example}
$(\mathbb{Z}, +, \cdot)$ is a commutative ring. The ring $M_n(\mathbb{R})$ of
$n \times n$ real matrices, with matrix addition and multiplication, is a
non-commutative ring for $n \ge 2$.
\end{example}

\section{Fields}

\begin{definition}[Field]
A \emph{field} is a commutative ring $(F, +, \cdot)$ in which $0 \neq 1$ and
every nonzero element has a multiplicative inverse: for every $a \in F$ with
$a \neq 0$ there exists $a^{-1} \in F$ such that $a \cdot a^{-1} = 1$.
\end{definition}

\begin{example}
$\mathbb{Q}$, $\mathbb{R}$, and $\mathbb{C}$ are fields. The ring
$\mathbb{Z}/p\mathbb{Z}$ is a field whenever $p$ is prime; for composite $n$
the ring $\mathbb{Z}/n\mathbb{Z}$ has zero divisors and is not a field.
\end{example}

\section{Why Fields Matter for ML}

A vector space is, by definition, an abelian group with scalar multiplication
by a field. Real-valued arithmetic in gradient descent, the covariance matrix
inverse in Gaussian likelihoods, and the advantage normalisation
$\hat{A}_i = (r_i - \mu)/\sigma$ in GRPO all silently use that $\mathbb{R}$ is
a field --- specifically, that $1/\sigma$ exists whenever $\sigma \neq 0$.

\section*{Exercises}
\begin{enumerate}
  \item Prove that $(a \cdot b)^{-1} = b^{-1} \cdot a^{-1}$ in any group.
  \item Find a nonzero element of $\mathbb{Z}/4\mathbb{Z}$ with no
        multiplicative inverse, and conclude that $\mathbb{Z}/4\mathbb{Z}$
        is not a field.
  \item Show that in any ring $R$, $0 \cdot a = 0$ for every $a \in R$.
\end{enumerate}

\end{document}
